\chapter{电磁感应}

电和磁并不是彼此孤立的两门学问。静止电荷产生电场，稳恒电流产生磁场，而一旦磁场随时间变化，或者导体在磁场中运动，电路中又会出现新的电动势与电流。这种由磁引电的现象就叫作\emph{电磁感应}。从微积分的角度看，电磁感应的核心只有一句话：\textbf{磁通量随时间的变化率决定感应电动势的大小，而负号决定感应效应总要反抗原来的变化。}

本章将用面积分、线积分与常微分方程，把高中阶段关于磁通量、法拉第定律、楞次定律、动生电动势、感生电动势、自感互感以及 RL 电路串成一个统一的理论链条。

\section{磁通量}

\begin{definition}
设空间中有磁场 $\vb{B}$，曲面 $S$ 上取有向面积元
\[
\dd \vb{A}=\uvec{n}\,\dd A,
\]
其中 $\uvec{n}$ 是曲面的单位法向量，则穿过曲面 $S$ 的\textbf{磁通量}定义为
\[
\Phi_B=\iint_S \vb{B}\cdot \dd \vb{A}=\iint_S \vb{B}\cdot \uvec{n}\,\dd A.
\]
\end{definition}

磁通量衡量的不是“有多少磁场”，而是“有多少磁场沿法向穿过该面”。因此，磁通量同时依赖于三个因素：磁场强弱、曲面面积大小、曲面的空间取向。

当磁场为匀强磁场，且曲面是平面时，上式化为熟悉的高中表达式
\[
\Phi_B=BA\cos\theta,
\]
其中 $\theta$ 是磁场方向与曲面法线之间的夹角，而不是与平面本身的夹角。

\begin{remark}
磁通量是一个\textbf{标量}，但它带正负号。正负号来自法向量的选取：如果 $\vb{B}$ 与 $\uvec{n}$ 大致同向，则 $\Phi_B>0$；若反向，则 $\Phi_B<0$。对闭合回路应用法拉第定律时，回路绕行方向与曲面法向必须满足右手定则，否则正负号会混乱。
\end{remark}

\begin{center}
\begin{tikzpicture}[scale=1.0]
  \coordinate (O) at (2.1,1.1);
  \coordinate (N) at (2.8,2.7);
  \coordinate (B) at (2.7,2.0);
  \fill[blue!8] (0,0) -- (3,0.8) -- (4.2,2.2) -- (1.2,1.4) -- cycle;
  \draw[thick] (0,0) -- (3,0.8) -- (4.2,2.2) -- (1.2,1.4) -- cycle;
  \draw[->, very thick, red!70!black] (O) -- (N) node[right] {$\uvec{n}$};
  \draw[->, very thick, blue!70!black] (0.6,2.7) -- (B) node[right] {$\vb{B}$};
  \path pic[draw, ->, "$\theta$", angle radius=0.55cm] {angle = B--O--N};
  \node at (2.1,-0.5) {曲面法向与磁场的夹角决定磁通量};
\end{tikzpicture}
\end{center}

若磁场或曲面并不均匀，就不能再用简单的 $BA\cos\theta$，而应回到面积分定义。把曲面分成许多小块，每一小块上的磁场近似不变，则
\[
\Phi_B\approx \sum_i \vb{B}_i\cdot \Delta \vb{A}_i,
\]
再令小块无限细分，就得到面积分极限。这正是“总量由局部量累加得到”的微积分思想。

\begin{example}
一面积为 $A$ 的矩形线圈置于匀强磁场 $\vb{B}$ 中，线圈法线与磁场夹角为 $\theta(t)=\omega t$。求任意时刻的磁通量，并指出磁通量变化最快的时刻。

\textbf{解：} 由定义
\[
\Phi_B(t)=BA\cos \theta(t)=BA\cos(\omega t).
\]
故其变化率为
\[
\dv{\Phi_B}{t}=-BA\omega \sin(\omega t).
\]
当 $\sin(\omega t)=\pm 1$ 时，磁通量变化率的绝对值最大，即
\[
\omega t=\frac{\pi}{2},\ \frac{3\pi}{2},\ \dots
\]
此时线圈平面与磁场平行，虽然瞬时磁通量本身为零，但磁通量变化得最快。
\end{example}

\begin{example}
半径为 $R$ 的圆面位于非均匀磁场中，磁场大小满足 $B(r)=B_0\left(1-\dfrac{r^2}{R^2}\right)$，方向处处垂直圆面向外。求穿过该圆面的磁通量。

\textbf{解：} 采用极坐标面积元 $\dd A=2\pi r\,\dd r$，因磁场与法线同向，故
\[
\Phi_B=\int_0^R B(r)\,2\pi r\,\dd r
=2\pi B_0 \int_0^R \left(1-\frac{r^2}{R^2}\right)r\,\dd r.
\]
计算得
\[
\Phi_B=2\pi B_0\left[\frac{r^2}{2}-\frac{r^4}{4R^2}\right]_0^R
=2\pi B_0\left(\frac{R^2}{2}-\frac{R^2}{4}\right)
=\frac{1}{2}\pi B_0 R^2.
\]
这个例子说明：即使圆面面积是 $\pi R^2$，有效磁通量也未必等于中心磁场乘面积，而必须对每一环带逐层积分。
\end{example}

\section{法拉第电磁感应定律}

实验表明，只要穿过闭合回路的磁通量发生变化，回路中就会出现感应电动势。变化可以来自磁场本身变化，也可以来自回路面积或方向变化。

\begin{law}[法拉第电磁感应定律]
对任意闭合回路 $C$ 及其所围曲面 $S$，感应电动势满足
\[
\mathcal{E}=-\dv{\Phi_B}{t}=-\dv{}{t}\iint_S \vb{B}\cdot \dd \vb{A}.
\]
若线圈有 $N$ 匝且每匝穿过的磁通量相同，则
\[
\mathcal{E}=-N\dv{\Phi_B}{t}.
\]
\end{law}

这里的负号不是数学装饰，而是物理内容的浓缩，它将在下一节由楞次定律解释。

从平均量的角度，若在时间间隔 $\Delta t$ 内磁通量改变了 $\Delta \Phi_B$，则平均感应电动势为
\[
\mathcal{E}_{\text{avg}}=-\frac{\Delta \Phi_B}{\Delta t}.
\]
令 $\Delta t\to 0$，就得到瞬时形式 $\mathcal{E}=-\dv{\Phi_B}{t}$。因此，法拉第定律是一个标准的导数定律：\textbf{电磁感应关心的不是磁通量有多大，而是它变化得有多快。}

\begin{remark}
很多题目中“磁通量为零”并不等于“没有感应电动势”。只要 $\dv{\Phi_B}{t}\neq 0$，即使某一时刻 $\Phi_B=0$，也仍然可能有感应电动势。上一节旋转线圈的例子就是如此。
\end{remark}

\begin{example}
一个 $N$ 匝线圈面积为 $A$，以角速度 $\omega$ 在匀强磁场 $B$ 中绕与磁场垂直的轴匀速转动。若初始时线圈法线与磁场同向，求感应电动势随时间的变化规律。

\textbf{解：} 每匝磁通量为
\[
\Phi_B(t)=BA\cos(\omega t).
\]
因此总磁链为 $N\Phi_B$，感应电动势
\[
\mathcal{E}(t)=-N\dv{\Phi_B}{t}=NBA\omega \sin(\omega t).
\]
于是交流发电机输出的是正弦交流电，其峰值为
\[
\mathcal{E}_{\max}=NBA\omega.
\]
这个公式说明，提高匝数、面积、磁场强度或转速，都能提高发电机输出电压。
\end{example}

\begin{figure}[htbp]
\centering
\begin{tikzpicture}
\begin{axis}[
  width=0.75\textwidth, height=0.5\textwidth,
  xlabel={时间 $t$}, ylabel={感应电动势 $\varepsilon/\varepsilon_0$},
  axis lines=middle,
  xmin=0, xmax=6.4, ymin=-1.1, ymax=1.1,
  grid=major, grid style={gray!30},
  thick,
]
\addplot[blue, domain=0:6.283, samples=180] {sin(deg(x))};
\end{axis}
\end{tikzpicture}
\caption{旋转线圈产生的感应电动势随时间作正弦变化}
\end{figure}

\begin{theorem}[广义法拉第定律]
若回路本身也在运动，则总电动势应写为
\[
\oint_C \left(\vb{E}+\vb{v}\times\vb{B}\right)\cdot \dd \vb{l}
=-\dv{}{t}\iint_S \vb{B}\cdot \dd \vb{A}.
\]
左边第一项对应感生电动势，第二项对应动生电动势。两者并不是彼此割裂的规律，而是同一电磁感应定律在不同情形下的表现。
\end{theorem}

这个定理告诉我们：当我们说“磁通量变化产生感应电动势”时，不必拘泥于变化究竟来自磁场、来自面积，还是来自回路位置。只要右边的磁通量作为时间函数发生改变，左边一定能在回路上找到与之对应的电动势来源。

\section{楞次定律}

法拉第定律给出大小，楞次定律给出方向。

\begin{law}[楞次定律]
感应电流（或感应电动势）所激发的磁场，总是\textbf{阻碍原来磁通量的变化}。
\end{law}

若穿过线圈的外磁通量增加，则感应电流产生的磁场要试图减小这种增加；若外磁通量减小，则感应电流产生的磁场要试图维持原来的磁通量。简言之，不是“反抗磁通量本身”，而是“反抗磁通量的变化”。

楞次定律实际上是能量守恒的必然结果。若感应电流总去“帮助”磁通量变化，那么系统便会无中生有地放大变化，产生免费的机械功或电功，这与能量守恒矛盾。负号 $-$ 正是在数学上排除了这种荒谬可能。

\begin{remark}
\textbf{方向示例：} 一根条形磁铁的 $N$ 极沿线圈轴线方向靠近静止线圈。此时穿过线圈的磁通量增大，根据楞次定律，线圈中感应电流产生的磁场必须阻碍这种“增大”，故线圈靠近磁铁的一侧应表现为 $N$ 极，从而排斥来近的磁铁。用右手螺旋定则判断，若从磁铁一侧观察，线圈中的感应电流应为逆时针方向。由于感应磁场对磁铁有排斥作用，外力若想继续把磁铁推近，就必须克服排斥力做正功；该机械功最终转化为电路中的焦耳热或磁场能量。
\end{remark}

\begin{remark}
判断感应电流方向时，最稳妥的步骤是：
\begin{enumerate}
  \item 先判断原磁通量是增大还是减小；
  \item 再判断感应磁场应当顺着原方向还是反着原方向；
  \item 最后用右手定则把感应磁场方向转化为电流方向。
\end{enumerate}
不要一上来就死记“靠近顺逆、远离顺逆”，因为一旦题目几何方向稍变，口诀就容易失效。
\end{remark}

\section{动生电动势}

\begin{definition}
当导体或回路的一部分在磁场中运动时，导体内自由电荷受到洛伦兹力
\[
\vb{F}_m=q\vb{v}\times\vb{B},
\]
从而沿导体定向分离，形成的电动势叫作\textbf{动生电动势}。其一般表达式为
\[
\mathcal{E}_{\text{mot}}=\int (\vb{v}\times\vb{B})\cdot \dd \vb{l}.
\]
\end{definition}

当导体棒长度为 $\ell$，速度 $\vb{v}$、磁场 $\vb{B}$ 与导体棒三者互相垂直时，上式简化为
\[
\mathcal{E}=B\ell v.
\]
这是高中题中最常见的切割磁力线模型。

\begin{center}
\begin{tikzpicture}[scale=1.0]
  \draw[thick] (0,0) -- (5,0);
  \draw[thick] (0,2.2) -- (5,2.2);
  \draw[very thick] (1.2,0) -- (2.0,2.2);
  \draw[->, thick] (2.2,1.1) -- (3.6,1.1) node[right] {$\vb{v}$};
  \foreach \x in {0.6,1.6,2.6,3.6,4.6} {
    \node at (\x,1.1) {$\otimes$};
  }
  \node at (4.4,1.75) {$\vb{B}$ 垂直纸面向里};
  \node[below] at (2.5,-0.1) {导轨};
  \node[above] at (2.6,2.3) {导轨};
  \node[left] at (1.2,1.1) {导体棒};
\end{tikzpicture}
\end{center}

在微观上，正电荷受到 $q\vb{v}\times\vb{B}$ 力而向导体一端聚集，负电荷向另一端聚集，直到静电力与磁力平衡。若导体两端再接成闭合回路，电荷便持续流动，形成感应电流。

\begin{example}
两平行光滑导轨间距为 $\ell$，置于竖直向下的匀强磁场 $B$ 中，一根金属棒以速度 $v$ 向右匀速滑动，导轨与棒构成闭合回路，总电阻为 $R$。求感应电流、安培力和外力功率。

\textbf{解：} 动生电动势大小为
\[
\mathcal{E}=B\ell v.
\]
故感应电流
\[
I=\frac{\mathcal{E}}{R}=\frac{B\ell v}{R}.
\]
棒中电流在磁场中受安培力，其大小为
\[
F=B I \ell=\frac{B^2\ell^2 v}{R}.
\]
根据楞次定律，该力方向必与运动方向相反，所以外力若要维持匀速，必须提供同样大小的拉力。外力功率为
\[
P_{\text{ext}}=Fv=\frac{B^2\ell^2 v^2}{R}.
\]
而电路中的焦耳热功率
\[
P_J=I^2R=\left(\frac{B\ell v}{R}\right)^2R=\frac{B^2\ell^2 v^2}{R}.
\]
两者相等，表明机械能通过电磁感应完整转化为内能。
\end{example}

\begin{remark}
“导体切割磁感线”只是动生电动势的一种形象说法。严格地说，关键不是是否看见“切割了线”，而是回路中电荷是否具有相对于磁场的速度，从而受到 $q\vb{v}\times\vb{B}$ 力。
\end{remark}

在导体整体进入磁场、并由涡流提供主要阻力的制动问题中，速度常近似满足指数衰减规律
\[
v(t)=v_0\mathrm{e}^{-t/\tau}.
\]
它表明制动力在速度较大时更明显，随后逐渐减弱。

\begin{figure}[htbp]
\centering
\begin{tikzpicture}
\begin{axis}[
  width=0.75\textwidth, height=0.5\textwidth,
  xlabel={时间 $t/\tau$}, ylabel={速度 $v/v_0$},
  axis lines=middle,
  xmin=0, xmax=5.2, ymin=0, ymax=1.08,
  grid=major, grid style={gray!30},
  thick,
]
\addplot[blue, domain=0:5, samples=180] {exp(-x)};
\end{axis}
\end{tikzpicture}
\caption{涡流制动中物体速度随时间逐渐衰减}
\end{figure}

\section{感生电动势}

当回路静止而磁场随时间变化时，电荷即使没有整体运动速度，仍然可能在空间中受到电场作用而沿回路运动。这种由变化磁场激发出的电动势叫作\textbf{感生电动势}。

\begin{law}[变化磁场产生电场]
若闭合回路静止，则感应电动势满足
\[
\oint_C \vb{E}\cdot \dd \vb{l}=-\dv{\Phi_B}{t}.
\]
这说明变化的磁场会在空间中激发环形电场。
\end{law}

与静电场不同，这种电场的电场线不是从正电荷出发到负电荷终止，而是绕成闭合曲线，因此它是\textbf{非保守电场}。用麦克斯韦方程的微分形式表示，即
\[
\curl \vb{E}=-\pdv{\vb{B}}{t}.
\]
右边不为零，说明此时电场具有旋度，不可能由单值静电势在全空间中简单描述。

\begin{remark}
动生电动势和感生电动势在本质上都属于电磁感应，但受力机制不同：
\begin{itemize}
  \item 动生电动势：电荷主要受磁场力 $q\vb{v}\times\vb{B}$；
  \item 感生电动势：电荷主要受变化磁场激发的感应电场力 $q\vb{E}$。
\end{itemize}
在复杂情形中，两种贡献可以同时存在，必须用广义法拉第定律统一处理。
\end{remark}

\begin{example}
半径为 $R$ 的圆形区域内存在均匀磁场 $\vb{B}(t)=B(t)\uvec{z}$，其外部磁场可忽略。已知 $\dv*{B}{t}=k$ 为常数。求以该区域圆心为中心、半径为 $r$ 的圆周上的感应电场大小 $E(r)$。

\textbf{解：} 由对称性，感应电场沿圆周切向分布，且同一圆周上大小相等，因此
\[
\oint \vb{E}\cdot \dd \vb{l}=E(r)\cdot 2\pi r.
\]
分两种情况：

\textbf{(1) $r<R$：} 取半径为 $r$ 的回路，穿过回路的磁通量为
\[
\Phi_B=B(t)\pi r^2,
\]
故
\[
E(r)\,2\pi r=-\dv{}{t}(B\pi r^2)=-\pi r^2 k.
\]
于是
\[
E(r)=-\frac{k r}{2}.
\]
其大小为
\[
|E(r)|=\frac{|k|r}{2}.
\]

\textbf{(2) $r>R$：} 真正有磁通变化的区域仍只有半径 $R$ 的圆盘，因此
\[
\Phi_B=B(t)\pi R^2,
\]
故
\[
E(r)\,2\pi r=-\pi R^2 k,
\qquad
|E(r)|=\frac{|k|R^2}{2r}.
\]

因此感应电场先随 $r$ 线性增大，超过磁场区域后再按 $1/r$ 衰减。这类分段函数正是面积分与线积分共同决定的结果。
\end{example}

\section{自感与互感}

电流变化会改变自身产生的磁场，因此不仅外磁场变化能引起感应电动势，回路中\textbf{自身电流的变化}也会导致感应现象。

\begin{definition}
若一个线圈中电流为 $I$，穿过每匝线圈的磁通量为 $\Phi_B$，且在线性介质中有
\[
N\Phi_B=LI,
\]
则称 $L$ 为该线圈的\textbf{自感系数}，简称\textbf{电感}。
\end{definition}

于是线圈中的自感电动势满足
\[
\mathcal{E}_L=-L\dv{I}{t}.
\]
电流增加时，自感电动势阻碍电流增加；电流减小时，自感电动势又阻碍电流减小。这正是楞次定律在“回路对自身变化”的体现。

对长直螺线管，若长度为 $l$，截面积为 $A$，单位长度匝数为 $n$，则其内部磁场近似为
\[
B=\mu_0 n I.
\]
总匝数 $N=nl$，于是磁链为
\[
N\Phi_B=N(BA)=nl(\mu_0 n I A).
\]
故电感近似为
\[
L=\mu_0 n^2 A l.
\]
这说明电感与线圈几何结构密切相关：匝数密、截面大、长度长，都有利于储存更多磁场能量。

\begin{definition}
若线圈 1 中的电流 $I_1$ 在线圈 2 中产生磁通量 $\Phi_{21}$，并满足
\[
N_2\Phi_{21}=M I_1,
\]
则称 $M$ 为两线圈之间的\textbf{互感系数}。线圈 2 中产生的感应电动势为
\[
\mathcal{E}_2=-M\dv{I_1}{t}.
\]
\end{definition}

互感是变压器工作的基础。交流电通过原线圈时，原线圈中的电流不断变化，于是铁芯中的磁通量不断变化，再在副线圈中激发感应电动势。

\begin{figure}[htbp]
\centering
\begin{tikzpicture}
\begin{axis}[
  width=0.75\textwidth, height=0.5\textwidth,
  xlabel={原线圈电流 $I_1$}, ylabel={磁链 $N_2\Phi_{21}$},
  axis lines=middle,
  xmin=0, xmax=5.2, ymin=0, ymax=4.5,
  grid=major, grid style={gray!30},
  thick,
]
\addplot[blue, domain=0:5, samples=120] {0.8*x};
\end{axis}
\end{tikzpicture}
\caption{在线性范围内，互感磁链与原线圈电流成正比}
\end{figure}

\begin{law}[理想变压器关系]
忽略漏磁、线圈电阻和能量损失时，理想变压器满足
\[
\frac{V_1}{V_2}=\frac{N_1}{N_2},
\qquad
\frac{I_1}{I_2}=\frac{N_2}{N_1},
\qquad
P_1\approx P_2.
\]
\end{law}

\begin{center}
\begin{tikzpicture}[scale=1.0]
  \draw[thick] (0,0) arc[start angle=-90,end angle=90,radius=0.25]
               (0.5,0) arc[start angle=-90,end angle=90,radius=0.25]
               (1.0,0) arc[start angle=-90,end angle=90,radius=0.25]
               (1.5,0) arc[start angle=-90,end angle=90,radius=0.25];
  \draw[thick] (4,0) arc[start angle=90,end angle=-90,radius=0.25]
               (4.5,0) arc[start angle=90,end angle=-90,radius=0.25]
               (5.0,0) arc[start angle=90,end angle=-90,radius=0.25]
               (5.5,0) arc[start angle=90,end angle=-90,radius=0.25];
  \draw[very thick, fill=gray!20] (2.2,-1.7) rectangle (3.3,1.7);
  \node at (2.75,0) {铁芯};
  \draw[thick] (0,-1.0) -- (0,-1.8);
  \draw[thick] (1.75,-1.0) -- (1.75,-1.8);
  \draw[thick] (4,-1.0) -- (4,-1.8);
  \draw[thick] (5.75,-1.0) -- (5.75,-1.8);
  \node at (0.9,2.0) {原线圈 $N_1$};
  \node at (4.9,2.0) {副线圈 $N_2$};
  \draw[->, blue!70!black, thick] (1.8,1.2) .. controls (2.5,1.8) and (3.0,1.8) .. (4.0,1.2);
  \draw[->, blue!70!black, thick] (4.0,-1.2) .. controls (3.2,-1.8) and (2.6,-1.8) .. (1.8,-1.2);
  \node[blue!70!black] at (2.8,2.2) {$\Phi_B(t)$};
\end{tikzpicture}
\end{center}

\section{RL 电路与电磁能量}

电感元件接入电路后，电流一般不再“瞬间跳变”，而要满足微分方程。最基本的模型是串联的 RL 电路。

\begin{center}
\begin{circuitikz}[american]
  \draw (0,0)
    to[battery1,l=$\mathcal{E}$] (0,3)
    to[R,l=$R$] (3,3)
    to[L,l=$L$] (6,3)
    -- (6,0) -- (0,0);
\end{circuitikz}
\end{center}

由基尔霍夫回路定律可得
\[
\mathcal{E}-IR-L\dv{I}{t}=0,
\]
即
\[
L\dv{I}{t}+RI=\mathcal{E}.
\]
这是一个一阶线性常微分方程。若在 $t=0$ 接通电源且初始电流 $I(0)=0$，其解为
\[
I(t)=\frac{\mathcal{E}}{R}\left(1-e^{-t/\tau}\right),
\qquad
\tau=\frac{L}{R}.
\]
其中 $\tau$ 称为\textbf{时间常数}。它刻画电流建立的快慢：$\tau$ 越大，电流变化越慢。

若把电源移去，只剩电阻和电感构成回路，则满足
\[
L\dv{I}{t}+RI=0,
\]
解为
\[
I(t)=I_0 e^{-t/\tau}.
\]
可见电感并不会维持电流永不消失，而只是通过释放原先储存在磁场中的能量，使电流缓慢衰减。

电感中的磁场能量为
\[
U_B=\frac{1}{2}LI^2.
\]
对于磁场近似均匀分布的区域，还可写成能量密度形式
\[
 u_B=\frac{B^2}{2\mu_0}.
\]
它与静电场能量密度 $u_E=\dfrac{1}{2}\varepsilon_0 E^2$ 共同说明：\textbf{电磁场本身就能储存能量}。

\begin{example}
一串联 RL 电路中，电源电动势为 $\mathcal{E}$，电阻为 $R$，电感为 $L$。求接通后 $t=\tau=L/R$ 时的电流，并求此时电感中储存的磁场能量占稳定时最大值的比例。

\textbf{解：} 电流随时间变化为
\[
I(t)=\frac{\mathcal{E}}{R}\left(1-e^{-t/\tau}\right).
\]
故当 $t=\tau$ 时，
\[
I(\tau)=\frac{\mathcal{E}}{R}\left(1-e^{-1}\right).
\]
稳定后最大电流为 $I_{\infty}=\dfrac{\mathcal{E}}{R}$，故最大磁场能量
\[
U_{\max}=\frac{1}{2}L\left(\frac{\mathcal{E}}{R}\right)^2.
\]
而 $t=\tau$ 时
\[
U(\tau)=\frac{1}{2}L I^2(\tau)=\frac{1}{2}L\left(\frac{\mathcal{E}}{R}\right)^2\left(1-e^{-1}\right)^2.
\]
因此所占比例为
\[
\frac{U(\tau)}{U_{\max}}=\left(1-e^{-1}\right)^2\approx 0.40.
\]
所以经过一个时间常数，电流已达到稳态值的约 $63\%$，但能量只达到约 $40\%$，因为能量与电流的平方成正比。
\end{example}

\begin{remark}
由 $U_B=\dfrac{1}{2}LI^2$ 可见，电感类似于力学中的“惯性”元件：它反抗电流变化，并把外界做功暂时存入磁场，再在需要时释放出来。这也是各种滤波、电能转换与暂态控制电路的基础。
\end{remark}

下面选取三道典型经典题型，展示如何把“公式套用”提升为“函数、导数与积分”的统一理解。

\section*{本章小结}

\begin{itemize}
  \item 磁通量是磁场对有向面积的面积分：$\Phi_B=\iint_S \vb{B}\cdot \dd \vb{A}$。
  \item 法拉第定律给出大小：$\mathcal{E}=-\dv{\Phi_B}{t}$；楞次定律给出方向：感应效应总阻碍原磁通量变化。
  \item 动生电动势来自 $q\vb{v}\times\vb{B}$，感生电动势来自变化磁场激发的环形电场。
  \item 自感与互感把“电流变化”转化为“感应电动势”，变压器则是互感的直接应用。
  \item RL 电路满足一阶线性微分方程，电感中的能量为 $U_B=\dfrac{1}{2}LI^2$，磁场能量密度为 $u_B=\dfrac{B^2}{2\mu_0}$。
\end{itemize}

